\documentclass[12pt,a4paper]{article}
\usepackage[utf8]{inputenc}
\usepackage[T1]{fontenc}
\usepackage{amsmath, amsthm, amssymb, amsfonts}
\usepackage{geometry}
\geometry{left=2.5cm, right=2.5cm, top=3cm, bottom=3cm}
\usepackage{hyperref}
\usepackage{setspace}
\setstretch{1.2}

\title{An Analytic and Algebraic Deconstruction of the Erd\H{o}s-Straus Conjecture: \\ Toward Symbolic Autoformalization}
\author{}
\date{}

\newtheorem{theorem}{Theorem}[section]
\newtheorem{lemma}[theorem]{Lemma}
\newtheorem{conjecture}[theorem]{Conjecture}
\newtheorem{definition}[theorem]{Definition}
\newtheorem{remark}[theorem]{Remark}
\newtheorem{proposition}[theorem]{Proposition}

\begin{document}
\maketitle

\begin{abstract}
The Erd\H{o}s-Straus conjecture postulates that for every integer $n \geq 2$, the rational number $4/n$ can be expressed as the sum of three positive unit fractions. This manuscript presents a rigorous, axiomatically grounded deconstruction of this conjecture. By synthesizing techniques from algebraic geometry, modular arithmetic, and combinatorial number theory, we isolate distinct congruence classes and provide complete demonstrations of polynomial identities bounding the solutions. Furthermore, we establish a rigid architectural framework intended for direct translation into the Lean 4 interactive theorem prover, thereby laying the groundwork for a fully verified symbolic resolution of the partial results.
\end{abstract}

\tableofcontents
\newpage

\section{Introduction and Axiomatic Foundations}

The study of Egyptian fractions, representations of rational numbers as sums of distinct unit fractions, dates back to antiquity. Among the modern inquiries into this domain, the Erd\H{o}s-Straus conjecture stands out for its structural simplicity and profound resistance to unconditional resolution. Formulated by Paul Erd\H{o}s and Ernst G. Straus in 1948, the conjecture asserts a specific upper bound on the number of terms required to represent fractions of the form $4/n$.

\begin{definition}[Erd\H{o}s-Straus Representation]
Let $n \in \mathbb{N}$ such that $n \geq 2$. An Erd\H{o}s-Straus representation of $n$ is an ordered triple $(x, y, z) \in \mathbb{N}^3$ such that $x > 0, y > 0, z > 0$ and the following relation holds:
\begin{equation}
\frac{4}{n} = \frac{1}{x} + \frac{1}{y} + \frac{1}{z}
\end{equation}
\end{definition}

\begin{conjecture}[Erd\H{o}s-Straus Conjecture]
\label{conj:es}
For all $n \in \mathbb{N}$ with $n \geq 2$, the set of Erd\H{o}s-Straus representations of $n$ is non-empty. In formal notation:
\begin{equation}
\forall n \in \mathbb{N}_{\geq 2}, \exists (x, y, z) \in \mathbb{Z}_{>0}^3 \text{ such that } 4xyz = n(xy + yz + zx)
\end{equation}
\end{conjecture}

To study this equation systematically, we clear denominators to obtain the equivalent Diophantine equation over the integers. Let us formally define the algebraic variety associated with the conjecture.

\begin{definition}[Associated Algebraic Variety]
For a fixed integer $n \geq 2$, the Erd\H{o}s-Straus variety $V_n \subset \mathbb{A}^3$ is defined as the zero locus of the polynomial $P_n(x,y,z) \in \mathbb{Z}[x,y,z]$:
\begin{equation}
P_n(x, y, z) = 4xyz - nxy - nyz - nzx
\end{equation}
The conjecture is equivalent to the statement that for all $n \geq 2$, the intersection $V_n \cap (\mathbb{Z}_{>0})^3$ is non-empty.
\end{definition}

We proceed by categorizing the domain $\mathbb{N}_{\geq 2}$ into distinct arithmetic progressions. By the Fundamental Theorem of Arithmetic, it suffices to prove the conjecture for prime numbers $n = p \geq 3$, since if $n = p \cdot k$, a representation $4/p = 1/x + 1/y + 1/z$ implies $4/n = 1/(xk) + 1/(yk) + 1/(zk)$.

\newpage
\section{Contextual Literature and Analogies}

The methodology for addressing such Diophantine problems typically relies on the construction of parametric polynomial identities. Mordell (1969) established the utility of such identities in showing that certain congruence classes admit solutions unconditionally.

\subsection{Analogous Frameworks: The Hasse Principle and Sieve Methods}

The problem at hand bears structural resemblance to the Waring's problem for cubes and the Goldbach conjecture, in the sense that local solutions (modulo primes) must be patched together to infer global integer solutions. However, the Hasse principle fails for the Erd\H{o}s-Straus equation in general due to the lack of sufficient local-global obstruction tools for non-homogeneous cubic surfaces.

Instead, the closest theoretical analog is the application of the large sieve method (as developed by Linnik and R\'{e}nyi, and utilized by Vaughan). Schinzel (1956) and Webb (1970) utilized probabilistic heuristics to estimate the asymptotic density of primes for which the conjecture holds. Elsholtz and Tao (2013) provided sharp upper bounds on the number of solutions, linking the problem to the distribution of divisors and primes in arithmetic progressions.

Our approach diverges from pure analytic density by returning to the algebraic roots: constructing explicit, rigorous identities for sub-classes of primes.

\newpage
\section{Strategy of Proof and Isolation of Lemmas}

We decompose the resolution of specific congruence classes into independent lemmas. The strategy relies on defining auxiliary parameters $(a, b, c, k) \in \mathbb{Z}^4$ to linearize the constraint equation.

By enforcing the structural form:
\begin{equation}
x = nab, \quad y = nac, \quad z = bck
\end{equation}
we can rewrite the primary equation $4xyz = nxy + nyz + nzx$. Let us substitute these generic forms into the target equation.

\begin{lemma}[Parametric Reduction Lemma]
\label{lem:parametric_reduction}
Let $n \in \mathbb{N}_{\geq 2}$. If there exist strictly positive integers $a, b, c, k$ such that:
\begin{equation}
4abc = abn + c(a+b)k
\end{equation}
then the triple $(x, y, z) = (nab, nac, bck)$ constitutes a valid Erd\H{o}s-Straus representation for $n$.
\end{lemma}

We will formulate two primary lemmas for specific modulo classes and provide their full demonstrations.
\begin{itemize}
\item \textbf{Lemma \ref{lem:mod4_case_1}}: The conjecture holds for all $n = 2 \pmod 3$.
\item \textbf{Lemma \ref{lem:mod4_case_2}}: The conjecture holds for all $n = 3 \pmod 4$.
\end{itemize}

\newpage
\section{Formal Demonstrations}

We now provide the complete, step-by-step rigorous demonstrations of the core lemmas. No algebraic manipulation is assumed or omitted.

\subsection{Proof of the Conjecture for $n = 3 \pmod 4$}
\begin{lemma}
\label{lem:mod4_case_2}
For any integer $n \geq 2$ such that $n = 3 \pmod 4$, there exist positive integers $x, y, z$ such that $4/n = 1/x + 1/y + 1/z$.
\end{lemma}

\begin{proof}
Let $n \in \mathbb{N}_{\geq 2}$. Assume that $n = 3 \pmod 4$.
By the definition of congruence modulo 4, there exists an integer $m \in \mathbb{Z}_{\geq 0}$ such that:
\begin{equation}
n = 4m + 3
\end{equation}

Let $k = m + 1$. Since $m \geq 0$, $k \geq 1$.
We have $4k - 1 = 4(m + 1) - 1 = 4m + 4 - 1 = 4m + 3 = n$.
Therefore, $n = 4k - 1$ for some integer $k \geq 1$.

We construct explicit integers $x, y, z$ as functions of $k$. Let us define:
\begin{equation}
x = k+1
\end{equation}
\begin{equation}
y = k(k+1)
\end{equation}
\begin{equation}
z = k(4k-1) = kn
\end{equation}
We now verify the properties of this representation strictly.

Positivity verification:
Since $k \geq 1$:
\begin{equation}
x = k + 1 \geq 2 > 0
\end{equation}
\begin{equation}
y = k(k + 1) \geq 1 \cdot 2 = 2 > 0
\end{equation}
\begin{equation}
z = k(4k - 1) = kn \geq 1 \cdot 3 = 3 > 0
\end{equation}

Reciprocal sum verification:
\begin{equation}
S = \frac{1}{x} + \frac{1}{y} + \frac{1}{z}
\end{equation}
\begin{equation}
S = \frac{1}{k+1} + \frac{1}{k(k+1)} + \frac{1}{kn}
\end{equation}
We combine the first two terms. The common denominator is $k(k+1)$:
\begin{equation}
\frac{1}{k+1} + \frac{1}{k(k+1)} = \frac{k}{k(k+1)} + \frac{1}{k(k+1)} = \frac{k+1}{k(k+1)}
\end{equation}
Since $k+1 > 0$, we can divide numerator and denominator by $k+1$:
\begin{equation}
\frac{k+1}{k(k+1)} = \frac{1}{k}
\end{equation}
Substituting this back into the sum $S$:
\begin{equation}
S = \frac{1}{k} + \frac{1}{kn}
\end{equation}
The common denominator is $kn$.
\begin{equation}
S = \frac{n}{kn} + \frac{1}{kn} = \frac{n+1}{kn}
\end{equation}
Recall that $n = 4k - 1$. Substitute $n$ in the numerator:
\begin{equation}
n + 1 = (4k - 1) + 1 = 4k
\end{equation}
Therefore:
\begin{equation}
S = \frac{4k}{kn}
\end{equation}
Since $k \geq 1$, we can divide out $k$:
\begin{equation}
S = \frac{4}{n}
\end{equation}
This concludes the rigorous demonstration for all $n = 3 \pmod 4$.
\end{proof}

\newpage
\subsection{Proof of the Conjecture for $n = 2 \pmod 3$}
\begin{lemma}
\label{lem:mod3_case_2}
For any integer $n \geq 2$ such that $n = 2 \pmod 3$, there exist positive integers $x, y, z$ such that $4/n = 1/x + 1/y + 1/z$.
\end{lemma}

\begin{proof}
Let $n \in \mathbb{N}_{\geq 2}$. Assume that $n = 2 \pmod 3$.
This implies $n$ can be written in the form $n = 3k - 1$ for some integer $k \ge 1$.
To establish the representation $4/n = 1/x + 1/y + 1/z$, we propose the following explicit parameterization:
\begin{equation}
x = k
\end{equation}
\begin{equation}
y = n
\end{equation}
\begin{equation}
z = kn
\end{equation}
Let us strictly verify positivity.
Since $n \geq 2$ and $n = 2 \pmod 3$, $n = 3k - 1 \ge 2 \implies 3k \ge 3 \implies k \ge 1$.
Therefore:
\begin{equation}
x = k \geq 1 > 0
\end{equation}
\begin{equation}
y = n \geq 2 > 0
\end{equation}
\begin{equation}
z = kn \geq 1 \cdot 2 = 2 > 0
\end{equation}

Reciprocal sum verification:
\begin{equation}
S = \frac{1}{x} + \frac{1}{y} + \frac{1}{z}
\end{equation}
Substitute the definitions:
\begin{equation}
S = \frac{1}{k} + \frac{1}{n} + \frac{1}{kn}
\end{equation}
The common denominator is $kn$.
\begin{equation}
S = \frac{n}{kn} + \frac{k}{kn} + \frac{1}{kn}
\end{equation}
\begin{equation}
S = \frac{n + k + 1}{kn}
\end{equation}
We must show this equals $4/n$, which is equivalent to showing $n + k + 1 = 4k$.
Recall the hypothesis $n = 3k - 1$. We substitute this into the numerator:
\begin{equation}
\text{Numerator} = (3k - 1) + k + 1
\end{equation}
\begin{equation}
\text{Numerator} = 3k + k - 1 + 1
\end{equation}
\begin{equation}
\text{Numerator} = 4k
\end{equation}
Thus, the sum evaluates to:
\begin{equation}
S = \frac{4k}{kn}
\end{equation}
Since $k \geq 1$, $k \neq 0$, we divide numerator and denominator by $k$:
\begin{equation}
S = \frac{4}{n}
\end{equation}
This concludes the rigorous demonstration for all $n = 2 \pmod 3$.
\end{proof}

\newpage
\subsection{Extensions to Higher Congruences: $n = 5 \pmod 8$}
\begin{lemma}
For any integer $n \geq 2$ such that $n = 5 \pmod 8$, there exist positive integers $x, y, z$ such that $4/n = 1/x + 1/y + 1/z$.
\end{lemma}
\begin{proof}
We assume $n = 5 \pmod 8$. Thus, $n = 8k - 3$ for some integer $k \ge 1$.
We utilize a more complex parameterization scheme.
Consider the algebraic identity:
\begin{equation}
\frac{4}{8k-3} = \frac{1}{2k} + \frac{1}{2k(8k-3)} + \frac{1}{k(8k-3)}
\end{equation}

Let us compute the sum of the first two terms:
\begin{equation}
\frac{1}{2k} + \frac{1}{2k(8k-3)} = \frac{8k-3+1}{2k(8k-3)} = \frac{8k-2}{2k(8k-3)} = \frac{4k-1}{k(8k-3)} = \frac{4k-1}{kn}
\end{equation}
We want the total sum to be $4/n = 4k/kn$.
Thus, the remainder needed is:
\begin{equation}
\frac{4k}{kn} - \frac{4k-1}{kn} = \frac{1}{kn}
\end{equation}
This perfectly generates the third term. Let us formally define our variables:
\begin{equation}
x = 2k
\end{equation}
\begin{equation}
y = kn
\end{equation}
\begin{equation}
z = 2kn
\end{equation}
Let us rigorously verify this proposed solution.
Positivity check: Since $n \geq 2$ and $n = 5 \pmod 8$, $n = 8k - 3 \ge 5 \implies 8k \ge 8 \implies k \ge 1$.
\begin{equation}
x = 2k \geq 2 > 0
\end{equation}
\begin{equation}
y = kn \geq 1 \cdot 5 = 5 > 0
\end{equation}
\begin{equation}
z = 2kn \geq 2 \cdot 5 = 10 > 0
\end{equation}
Sum of reciprocals:
\begin{equation}
S = \frac{1}{x} + \frac{1}{y} + \frac{1}{z}
\end{equation}
\begin{equation}
S = \frac{1}{2k} + \frac{1}{kn} + \frac{1}{2kn}
\end{equation}
The common denominator is $2kn$.
\begin{equation}
S = \frac{n}{2kn} + \frac{2}{2kn} + \frac{1}{2kn}
\end{equation}
\begin{equation}
S = \frac{n + 2 + 1}{2kn}
\end{equation}
\begin{equation}
S = \frac{n + 3}{2kn}
\end{equation}
Substitute $n = 8k - 3$:
\begin{equation}
\text{Numerator} = (8k - 3) + 3 = 8k
\end{equation}
Therefore:
\begin{equation}
S = \frac{8k}{2kn} = \frac{4}{n}
\end{equation}
This concludes the rigorous demonstration for all $n = 5 \pmod 8$.
\end{proof}

\newpage
\section{The Modulo 24 Sub-cases Framework}

It is a well known result (due to Mordell) that the Erd\H{o}s-Straus conjecture is true for all $n$ except possibly those $n = 1, 121, 169, 289, 361, 529 \pmod{840}$. To illustrate the method of covering congruences, we give the explicit identities for all non-trivial congruences modulo 24.
Prime numbers modulo 24 can only take the values $1, 5, 7, 11, 13, 17, 19, 23$. We have already solved $n = -1 \pmod 4$ (which covers $3, 7, 11, 15, 19, 23$).
We have already solved $n = -1 \pmod 3$ (which covers $2, 5, 8, 11, 14, 17, 20, 23$).
We have solved $n = 5 \pmod 8$ (which covers $5, 13, 21$).

Thus, the only remaining prime classes modulo 24 to check are:
- $n = 1 \pmod{24}$

Observe that:
$n = 13 \pmod{24}$ implies $n = 5 \pmod 8$, which is covered.
$n = 17 \pmod{24}$ implies $n = 2 \pmod 3$, which is covered.
$n = 19 \pmod{24}$ implies $n = 3 \pmod 4$, which is covered.
$n = 23 \pmod{24}$ implies $n = 3 \pmod 4$ and $n = 2 \pmod 3$, covered.

Therefore, the only remaining unresolved class modulo 24 is $n = 1 \pmod{24}$.
This demonstrates the extraordinary power of these elementary polynomial identities to cover an asymptotic density approaching 1.

\newpage
\section{Further Expansions and Asymptotic Properties}

To fully comprehend the structural behavior of the Erd\H{o}s-Straus conjecture, it is essential to explore the density of exceptions in broader moduli bases. Although polynomial identities cover an arbitrary fraction of the integer space, the set of uncovered primes for any finite base of identities remains infinite. This section develops the analytical framework used to bound the exceptional set.

Let $E(X)$ denote the number of integers $n \le X$ for which the Erd\H{o}s-Straus conjecture is false. Vaughan (1970) proved that there exists a positive constant $c$ such that:
\begin{equation}
E(X) \le X \exp\left( -c \frac{\log X}{\log \log X} \right)
\end{equation}

By extending the identities established in the previous sections, one can derive more refined congruences. For example, the class $n \equiv 1 \pmod{24}$ can be further partitioned using modulo 840, which is $8 \times 3 \times 5 \times 7$.

If $n \equiv 1 \pmod{24}$, we consider its residue modulo 5. The cases $n \equiv 2, 3, 4 \pmod 5$ admit direct identities similar to those demonstrated for modulo 3 and 4. The only problematic residue is $n \equiv 1 \pmod 5$. By aggregating these properties across multiple primes, one obtains the celebrated result that the only primes unconditionally resisting polynomial parameterization are those congruent to squares modulo 840.

\newpage
\section{Algorithmic Verification and Computational Bounds}

In parallel to analytical methods, computational verifications play a crucial role in validating the conjecture up to vast numerical limits. As of current literature, the conjecture has been verified for all $n \le 10^{17}$.

The algorithm employed for these verifications is heavily optimized by exploiting the very polynomial identities we have proven. A naive search for $x, y, z$ takes $O(n^2)$ time. However, by sieving through known congruence classes, the algorithm only tests $n$ which are not covered by any known identity.

For a given $n$, the algorithm clears the denominator to form $4xyz = nxy + nyz + nzx$, and then iterates through divisors of the quadratic forms associated with the equation. Let $A = 4x - n$. The equation can be rewritten as:
\begin{equation}
4xyz = nxy + n(x+y)z \implies z(4xy - n(x+y)) = nxy
\end{equation}
By setting $x = \lceil n/4 \rceil + i$, and bounding the search space of $i$, computational systems can efficiently locate a solution or prove its absence within the local search domain.

\newpage
\section{Generalization: The Sierpiński Conjecture}

A natural extension of the Erd\H{o}s-Straus conjecture is the Sierpiński conjecture, which posits that for any integer $n \ge 4$, the equation
\begin{equation}
\frac{5}{n} = \frac{1}{x} + \frac{1}{y} + \frac{1}{z}
\end{equation}
has a solution in positive integers $x, y, z$.

The techniques developed for $4/n$ translate almost directly to $5/n$, albeit with a reduction in the density of polynomial identity coverage. For instance, the identity for $n \equiv 2 \pmod 3$ for $4/n$ can be structurally modified to cover specific classes for $5/n$.

Let $n = 3k - 1$. For $5/n$, we observe:
\begin{equation}
\frac{5}{3k-1} = \frac{3}{3k-1} + \frac{2}{3k-1}
\end{equation}
While not yielding a direct 3-term expansion in all cases, the systematic breakdown into fractional components allows for similar modular classification. The study of the $5/n$ case remains an active area of combinatorial number theory, mirroring the challenges of $4/n$ but with increased algebraic complexity.

\newpage
\section{Architecture for Autoformalization}

To transition from the human-readable proofs above to a fully verifiable symbolic system, we present the structural schema for Lean 4.
The axioms, types, and lemmas are strictly typed to guarantee compatibility with `Mathlib`.

\begin{verbatim}
import Mathlib.Data.Nat.Basic
import Mathlib.Tactic.Omega
import Mathlib.Tactic.Ring

-- Definition of the Erdős-Straus property for a given n
def ErdosStrausProperty (n : Nat) : Prop :=
  exists x y z : Nat, x > 0 /\ y > 0 /\ z > 0 /\ 4 * x * y * z = n * (x * y + y * z + z * x)

-- The main conjecture statement
def ErdosStrausConjecture : Prop :=
  forall n : Nat, n >= 2 -> ErdosStrausProperty n

-- Lemma for n = 3 mod 4
lemma erdos_straus_mod_4_case_3 (n k : Nat) (h_eq : n = 4 * k + 3) :
  ErdosStrausProperty n := by
  -- We reconstruct the substitution n = 4(k+1) - 1.
  -- Let m = k + 1, so n = 4m - 1.
  -- This requires precise handling in Lean.
  sorry
  -- Il s'agit d'une esquisse de preuve incomplete destinee a une autoformalisation future.

-- Lemma for n = 2 mod 3
lemma erdos_straus_mod_3_case_2 (n m : Nat) (h_eq : n + 1 = 3 * m) (h_m : m >= 1) :
  ErdosStrausProperty n := by
  -- The variables x, y, z constructed in the mathematical proof.
  use m
  use n
  use m * n
  -- Prove strict positivity
  have hx : m > 0 := by omega
  have hy : n > 0 := by omega
  have hz : m * n > 0 := by
    apply mul_pos hx hy
  refine ⟨hx, hy, hz, ?_⟩
  -- Prove the algebraic relation
  -- 4 * x * y * z = n * (x * y + y * z + z * x)
  -- 4 * m * n * (m * n) = n * (m * n + n * (m * n) + (m * n) * m)
  sorry
  -- Il s'agit d'une esquisse de preuve incomplete destinee a une autoformalisation future.

-- Lemma for n = 5 mod 8
lemma erdos_straus_mod_8_case_5 (n k : Nat) (h_eq : n + 3 = 8 * k) (h_k : k >= 1) :
  ErdosStrausProperty n := by
  use 2 * k
  use k * n
  use 2 * k * n
  sorry
  -- Il s'agit d'une esquisse de preuve incomplete destinee a une autoformalisation future.
\end{verbatim}

The strategic use of polynomial identities in Lean requires the `ring` tactic to automatically resolve the equality goals once the terms are instantiated. The primary challenge lies in the type conversions (e.g., operating in $\mathbb{Z}$ vs $\mathbb{N}$ due to subtraction, which is truncated in $\mathbb{N}$). Our parameterizations strictly avoid subtraction in the definitions of $x, y, z$, thereby remaining entirely within $\mathbb{N}$ and simplifying the autoformalization process.

\newpage
\section{Conclusion}

This document has provided an exhaustive, purely axiomatic derivation of the polynomial identities necessary to solve the Erd\H{o}s-Straus conjecture for multiple dense congruence classes. Every algebraic substitution, factorization, and common-denominator expansion has been explicitly documented.

The provided Lean 4 architectural skeleton translates these classical Diophantine techniques into a modern type-theoretic framework. The isolation of these lemmas demonstrates that while the global conjecture remains open, the structure of its partial solutions is deeply regular and amenable to automated symbolic verification.

\vspace{2cm}
\begin{center}
\textbf{End of Document}
\end{center}

\end{document}
